% answers_short.tex —— 2018 年·天津卷·文科 · 简短答案速查
% 由 tools/make_short_answers.py 自 answers.tex 的 \gksolhead/\gkanswer 提取，
% 并与 2018/2018_数学真题_天津卷文_答案要点.md 逐题比对后生成，勿手改（改源 MD 后重跑）。

\begin{gknotice}
\textbf{使用说明}：本节只给最终答案便于快速核对，完整推导见解析版。答案为编写者独立推导，\textbf{非官方标准答案}。
\end{gknotice}

\gkpart{选择题}
\begin{center}\begin{tabular}{ccccccccc}
\toprule
题号 & 1 & 2 & 3 & 4 & 5 & 6 & 7 & 8 \\
\midrule
答案 & C & C & A & B & D & A & A & C \\
\bottomrule
\end{tabular}\end{center}

\gkpart{填空题}
\begin{itemize}[leftmargin=2.6em]
  \item[\textbf{9.}] $4-\mathrm i$
  \item[\textbf{10.}] $\mathrm e$
  \item[\textbf{11.}] $1/3$
  \item[\textbf{12.}] $(x-1)^2+y^2=1$
  \item[\textbf{13.}] $1/4$
  \item[\textbf{14.}] $[1/8,2]$
\end{itemize}

\gkpart{解答题}
\begin{itemize}[leftmargin=2.6em]
  \item[\textbf{15.}] (1) 甲、乙、丙三个年级依次抽取 $3,2,2$ 人；(2) 抽取 2 名同学的所有可能结果共有 21 种；$P(M)=5/21$
  \item[\textbf{16.}] (1) $B=\pi/3$；(2) $b=\sqrt7$，$\sin(2A-B)=3\sqrt3/14$
  \item[\textbf{17.}] (1) 证明见解析版；(2) $\sqrt{13}/26$；(3) $\sqrt3/4$
  \item[\textbf{18.}] (1) $S_n=n(n+1)/2$，$T_n=2^n-1$；(2) $n=4$
  \item[\textbf{19.}] (1) $x^2/9+y^2/4=1$；(2) $k=-1/2$
  \item[\textbf{20.}] (1) $y=-x$；(2) 极大值为 $6\sqrt3$，极小值为 $-6\sqrt3$；(3) $(-\infty,-\sqrt{10})\cup(\sqrt{10},+\infty)$
\end{itemize}
